A kernel classifier that labels each flower by the class mean it is nearer to, measured in the kernel's feature space, is compared across four kernels (dot product, polynomial of degree 2 and 3, and Gaussian with $\sigma = 1/\sqrt{2}$) on four binary problems built from the Iris data. Each of five training shares is evaluated on \N{n_splits_ratio} random splits. On the three problems that are not linearly separable, the Gaussian kernel has the highest mean test accuracy, \N{gauss_p234}, against \N{dot_p234}, \N{d2_p234}, and \N{d3_p234} for the dot product and the degree-2 and degree-3 polynomials, and it is at least as accurate as the best other kernel on \N{gauss_ge_min}\% to \N{gauss_ge_max}\% of splits. When versicolor is separated from the other two species, the dot-product and polynomial kernels score below the majority-class rate (the larger class's share of each test set) on average, although in an added comparison with standardized features, the degree-2 kernel rises from \N{s_d2_p3_raw} to \N{s_d2_p3_std}. The training share, from 50\% to 90\%, does not change the observed mean accuracy averaged over all problem--kernel pairs by more than \N{all13_span}. This extended study uses raw measurements, one fixed Gaussian width, and one data set, and finds that the single split per training share in the original coursework submission gave higher accuracies on average than random splits do.
